% European Heart Journal HNMT Heart Failure 2026 教學講義
% Histamine N-methyltransferase 與心衰竭
\documentclass[12pt,a4paper]{article}
% Drake_preamble.tex - LaTeX Preamble Template
% 作者: 謝慕揚, MD, PhD, FESC
% 
% 使用說明:
% 1. 表格請使用 longtable，不要有垂直分隔線 |
% 2. 表格內換行使用 \newline，不要使用 \\
% 3. 藥物名稱、檢驗、檢查請維持英文
% 4. Reference 格式請使用 \href 加上驗證過的 DOI

% Including essential packages for Chinese typesetting
\usepackage{xeCJK}
\usepackage{fontspec}
% Configuring Chinese font with Noto Serif CJK TC, fallback to SimSun
\IfFontExistsTF{Noto Serif CJK TC}{
  \setCJKmainfont{Noto Serif CJK TC}
  \setCJKsansfont{Noto Serif CJK TC}
  \setCJKmonofont{Noto Serif CJK TC}
}{\setCJKmainfont{SimSun}
  \setCJKsansfont{SimSun}
  \setCJKmonofont{SimSun}
}

% Including other necessary packages
\usepackage{geometry}
\usepackage{titlesec}
\usepackage{enumitem}
\usepackage{xcolor}
\usepackage{colortbl}
\usepackage{tcolorbox}
\usepackage{booktabs}
\usepackage{longtable}
\usepackage{array}
\usepackage{graphicx}
\usepackage{tikz}
\usepackage{fancyhdr}
\usepackage{multicol}
\usepackage{datetime2}
\usepackage{hyperref}
\usepackage{mdframed}
\usepackage{amsmath}
\usepackage{amssymb}
\usepackage{multirow}

\hypersetup{
    colorlinks=true,
    linkcolor=emergency_blue,
    citecolor=emergency_blue,
    urlcolor=emergency_blue,
    pdfborder={0 0 0}
}

% Defining page geometry
\geometry{
  top=2.5cm,
  bottom=2.5cm,
  left=2cm,
  right=2cm,
  headheight=25pt
}

% Adjusting topmargin to compensate for increased headheight
\addtolength{\topmargin}{-10pt}

% Defining colors
\definecolor{emergency_red}{RGB}{186,24,27}
\definecolor{emergency_blue}{RGB}{0,114,188}
\definecolor{emergency_gray}{RGB}{120,120,120}
\definecolor{highlight_yellow}{RGB}{255,250,205}

% Setting title formats
\titleformat{\section}[block]{\Large\bfseries\color{emergency_red}}{\thesection}{10pt}{}
\titleformat{\subsection}[block]{\large\bfseries\color{emergency_blue}}{\thesubsection}{8pt}{}
\titleformat{\subsubsection}[block]{\normalsize\bfseries\color{emergency_gray}}{\thesubsubsection}{6pt}{}

% Defining custom environment for important notes
\newmdenv[
    linecolor=emergency_red,
    backgroundcolor=highlight_yellow,
    linewidth=2pt,
    topline=true,
    bottomline=true,
    leftline=true,
    rightline=true,
    innertopmargin=10pt,
    innerbottommargin=10pt,
    innerleftmargin=10pt,
    innerrightmargin=10pt
]{importantbox}

% Defining note command with tcolorbox
\newcommand{\note}[1]{
    \par
    \noindent
    \begin{tcolorbox}[
        colback=emergency_blue!5!white,
        colframe=emergency_blue,
        title=\textbf{重點提醒},
        fonttitle=\bfseries,
        boxrule=0.5pt,
        arc=3pt,
        left=6pt,
        right=6pt,
        top=6pt,
        bottom=6pt
    ]
    #1
    \end{tcolorbox}
}

% Key findings box
\newcommand{\keyfinding}[1]{
    \par
    \noindent
    \begin{tcolorbox}[
        colback=yellow!10!white,
        colframe=orange!75!black,
        title=\textbf{核心發現摘要},
        fonttitle=\bfseries,
        boxrule=0.5pt,
        arc=3pt
    ]
    #1
    \end{tcolorbox}
}

% Defining custom date format for ISO 8601
\DTMnewdatestyle{iso}{
  \renewcommand{\DTMdisplaydate}[4]{\number##1-\number##2-\number##3}
  \renewcommand{\DTMDisplaydate}[4]{\number##1-\number##2-\number##3}
}
\DTMsetdatestyle{iso}
\newcommand{\mydate}{\DTMtoday}

\usetikzlibrary{positioning}

% Custom header for this document
\pagestyle{fancy}
\fancyhf{}
\fancyhead[L]{\textcolor{emergency_red}{European Heart Journal 教學講義}}
\fancyhead[R]{\textcolor{emergency_blue}{HNMT 與心衰竭}}
\fancyfoot[C]{\textcolor{emergency_gray}{\thepage}}
\renewcommand{\headrulewidth}{1pt}
\renewcommand{\headrule}{\hbox to\headwidth{\color{emergency_red}\leaders\hrule height \headrulewidth\hfill}}

\begin{document}

%============================================================
% Title Page
%============================================================
\begin{center}
{\LARGE\bfseries\textcolor{emergency_red}{European Heart Journal}}\\[0.5cm]
{\Large\textbf{Histamine N-methyltransferase Upregulation,\\Cardiac Hypertrophy, and Heart Failure}}\\[0.3cm]
{\large Histamine N-methyltransferase 上調、心臟肥大與心衰竭}\\[0.5cm]
{\normalsize \href{https://doi.org/10.1093/eurheartj/ehaf1065}{Eur Heart J 2026. DOI: 10.1093/eurheartj/ehaf1065}}\\[0.3cm]
{\normalsize 整理：謝慕揚 MD, PhD, FESC \quad 日期：\mydate}
\end{center}

\vspace{0.5cm}

\keyfinding{
\textbf{核心發現：}
\begin{itemize}[noitemsep]
    \item HNMT（Histamine N-methyltransferase）在心衰竭患者心臟組織中表現上升
    \item HNMT 透過 SAM/EZH2/FZD2/CaMKII 軸促進心臟肥大與心衰竭
    \item Amodiaquine（HNMT 抑制劑）具有治療心衰竭的潛力
    \item 尿液 N-methylhistamine 可作為非侵入性心衰竭生物標記
\end{itemize}
}

%============================================================
\section{研究背景}
%============================================================

\begin{itemize}
    \item Histamine 訊號傳遞在心衰竭中扮演重要角色
    \item \textbf{HNMT}（Histamine N-methyltransferase）是心臟中清除 histamine 的主要酵素
    \item HNMT 在心衰竭中的角色尚未明確
    \item 本研究探討 HNMT 介導的代謝-表觀遺傳交互作用在心衰竭中的機轉
\end{itemize}

%============================================================
\section{研究方法}
%============================================================

\begin{longtable}{p{4cm}p{10cm}}
\toprule
\textbf{方法} & \textbf{內容} \\
\midrule
\endfirsthead
\bottomrule
\endfoot
動物模型 & • 心肌細胞特異性 Hnmt 基因剔除小鼠\newline • Hnmt 過度表現小鼠 \\
\midrule
心衰竭誘導 & Transverse aortic constriction（TAC）\\
\midrule
分子分析 & • RNA sequencing\newline • 標靶代謝體學 \\
\midrule
臨床驗證 & 53 位心衰竭患者 vs 55 位對照組\newline 測量尿液 N-methylhistamine 濃度 \\
\bottomrule
\end{longtable}

%============================================================
\section{主要發現}
%============================================================

\subsection{HNMT 在心衰竭中上調}

\begin{itemize}
    \item 心衰竭患者心臟組織：HNMT 表現 $\uparrow$
    \item TAC 手術小鼠心臟：HNMT 表現 $\uparrow$
    \item Phenylephrine 處理的新生小鼠心肌細胞：HNMT 表現 $\uparrow$
    \item 心衰竭患者尿液 N-methylhistamine $\uparrow$，且與心衰竭嚴重度正相關
\end{itemize}

\subsection{HNMT 基因操作的效果}

\begin{longtable}{p{5cm}p{9cm}}
\toprule
\textbf{操作} & \textbf{結果} \\
\midrule
\endfirsthead
\bottomrule
\endfoot
心肌細胞特異性 Hnmt 剔除 & 減輕 TAC 或 Angiotensin II 誘導的心臟肥大與心衰竭 \\
\midrule
Hnmt 過度表現 & 損害心臟功能 \\
\midrule
Amodiaquine（HNMT 抑制劑） & 改善 TAC 誘導的心臟功能障礙 \\
\bottomrule
\end{longtable}

%============================================================
\section{分子機轉}
%============================================================

\begin{center}
\begin{tikzpicture}[
    node distance=0.6cm,
    box/.style={rectangle, draw=emergency_blue, fill=emergency_blue!10,
                text width=3.5cm, minimum height=0.8cm, align=center,
                rounded corners, font=\small},
    arrow/.style={->, thick, >=stealth}
]
    \node[box, fill=emergency_red!20, draw=emergency_red] (hnmt) {HNMT $\uparrow$};
    \node[box, below=of hnmt] (sam) {SAM $\downarrow$};
    \node[box, below=of sam] (ezh2) {EZH2 功能受損};
    \node[box, below=of ezh2] (h3k27) {H3K27me3 $\downarrow$\\（Fzd2 啟動子）};
    \node[box, below=of h3k27] (fzd2) {FZD2 $\uparrow$};
    \node[box, below=of fzd2] (wnt) {WNT/Ca\textsuperscript{2+}/CaMKII\\pathway 活化};
    \node[box, below=of wnt, fill=emergency_red!20, draw=emergency_red] (hf) {心臟肥大 / 心衰竭};

    \draw[arrow] (hnmt) -- (sam);
    \draw[arrow] (sam) -- (ezh2);
    \draw[arrow] (ezh2) -- (h3k27);
    \draw[arrow] (h3k27) -- (fzd2);
    \draw[arrow] (fzd2) -- (wnt);
    \draw[arrow] (wnt) -- (hf);
\end{tikzpicture}
\end{center}

\subsection{機轉說明}

\begin{enumerate}
    \item \textbf{HNMT 上調}：心衰竭時 HNMT 表現增加
    \item \textbf{SAM 減少}：HNMT 消耗 S-adenosylmethionine（SAM）作為甲基供體
    \item \textbf{EZH2 功能受損}：SAM 是 EZH2（enhancer of zeste homologue 2）的輔因子
    \item \textbf{H3K27me3 降低}：EZH2 功能下降導致 Fzd2 啟動子的 H3K27me3 修飾減少
    \item \textbf{FZD2 上調}：表觀遺傳抑制解除，FZD2（Frizzled-2）表現上升
    \item \textbf{CaMKII 活化}：FZD2 活化 WNT/calcium/CaMKII 訊號路徑
    \item \textbf{心衰竭}：CaMKII 過度活化促進病理性心臟肥大與心衰竭
\end{enumerate}

%============================================================
\section{臨床意義}
%============================================================

\begin{importantbox}
\textbf{治療標靶：}
\begin{itemize}[noitemsep]
    \item \textbf{Amodiaquine}：已知的 HNMT 抑制劑，原為抗瘧疾藥物
    \item 動物實驗顯示可改善 TAC 誘導的心臟功能障礙
    \item 可能成為心衰竭的新治療策略
\end{itemize}

\vspace{0.3cm}

\textbf{生物標記：}
\begin{itemize}[noitemsep]
    \item 尿液 N-methylhistamine 可作為非侵入性心衰竭生物標記
    \item 濃度與心衰竭嚴重度正相關
    \item 優點：非侵入性、易於收集
\end{itemize}
\end{importantbox}

%============================================================
\section{結論}
%============================================================

\begin{enumerate}
    \item HNMT 透過 SAM/FZD2/CaMKII 軸加劇病理性心臟肥大與心衰竭
    \item Amodiaquine（HNMT 抑制劑）展現治療心衰竭的潛力
    \item 尿液 N-methylhistamine 可作為非侵入性心衰竭生物標記
\end{enumerate}

%============================================================
\section{參考文獻}
%============================================================

\begin{enumerate}
    \item Zhang J, et al. Histamine N-methyltransferase upregulation, cardiac hypertrophy, and heart failure. \href{https://doi.org/10.1093/eurheartj/ehaf1065}{\textit{Eur Heart J}. 2026. DOI: 10.1093/eurheartj/ehaf1065}
\end{enumerate}

\end{document}
